\documentclass[12pt,letterpaper]{exam}
\usepackage[lmargin=1in,rmargin=1in,tmargin=1in,bmargin=1in]{geometry}
\usepackage{../style/exams}

% -------------------
% Course & Exam Information
% -------------------
\newcommand{\course}{MATH 344: Exam 2}
\newcommand{\term}{Spring ---\textsubscript{\textsubscript{1}} 2026}
\newcommand{\examdate}{03/27/2026}
\newcommand{\timelimit}{50 Minutes}

\setbool{hideans}{false} % Student: True; Instructor: False

\usepackage{array}
\newcolumntype{C}[1]{>{\centering\arraybackslash}p{#1}}
\newcommand{\clinevsol}[1]{\ifbool{hideans}{}{#1}}
\newcommand{\vanbox}[1]{\ifbool{hideans}{#1}{\boxed{#1}}}
\DeclareMathOperator{\rank}{rank}
\DeclareMathOperator{\nullity}{nullity}
\DeclareMathOperator{\mspan}{span}

% -------------------
% Content
% -------------------
\begin{document}

\examtitle
\instructions{Write your name on the appropriate line on the exam cover sheet. This exam contains \numpages\ pages (including this cover page) and \numquestions\ questions. Check that you have every page of the exam. Answer the questions in the spaces provided on the question sheets. Be sure to answer every part of each question and show all your work. If you run out of room for an answer, continue on the back of the page --- being sure to indicate the problem number.} 
\scores
\bottomline
\newpage


% -------------------
% Questions
% -------------------
\begin{questions}

% Question 1
\newpage
\question[10] Let $\mathbf{b}= \begin{pmatrix} 3 \\ -4 \\ 11 \end{pmatrix}$ and $A$ be the matrix with LU factorization
	\[
	A= \begin{pmatrix} 1 & 0 & 0 \\ -1 & 1 & 0 \\ 4 & -3 & 1 \end{pmatrix}%
	\begin{pmatrix} 1 & 2 & -1 \\ 0 & 3 & 1 \\ 0 & 0 & 4 \end{pmatrix}
	\]

\begin{enumerate}[(a)]
\item Compute $\det A$. \pspace
\psol{\itshape We use the fact that $\det(MN)= \det M \cdot \det N$ and the determinant of a triangular matrix is the product of the diagonal entries. Therefore, we have\dots}
	\[
	\psol{\det A= \det(LM)= \det L \cdot \det M= 1 \cdot (1 \cdot 3 \cdot 4)= 12}
	\]

\item Use the LU factorization above to solve the matrix-vector equation $A \mathbf{x}= \mathbf{b}$. \pspace
\psol{\itshape We know that $A= LU$, so $\mathbf{b}= A\mathbf{x}= (LU) \mathbf{x}= L(U \mathbf{x})$. We also know that $U\mathbf{x}$ is a vector, call this vector $\mathbf{y}$. So, we have $L\mathbf{y}= \mathbf{b}$. Therefore, we have\dots}
	\[
	\psol{\begin{pmatrix} 1 & 0 & 0 \\ -1 & 1 & 0 \\ 4 & -3 & 1 \end{pmatrix} \begin{pmatrix} y_1 \\ y_2 \\ y_3 \end{pmatrix}= \begin{pmatrix} 3 \\ -4 \\ 11 \end{pmatrix}}
	\]
\psol{\itshape We solve this system using forward-substitution:}
	\[
	\begin{aligned}
	\psol{1y_1 + 0y_2 + 0y_3= 3} &\psol{\quad\Longrightarrow\quad y_1= 3} \\
	\psol{-1y_1 + 1y_2 + 0y_3= -4} &\psol{\quad\Longrightarrow\quad -3 + y_2= -4} &\psol{\quad\Longrightarrow\quad y_2= -1} \\
	\psol{4y_1 - 3y_2 + 1y_3= 11} &\psol{\quad\Longrightarrow\quad 4(3) - 3(-1) + y_3= 11} &\psol{\quad\Longrightarrow\quad y_3= -4}
	\end{aligned}
	\]
\psol{\itshape Because $U\mathbf{x}= \mathbf{y}$ and $\mathbf{y}= \begin{pmatrix} 3 \\ -1 \\ -4 \end{pmatrix}$, we know\dots}
	\[
	\psol{\begin{pmatrix} 1 & 2 & -1 \\ 0 & 3 & 1 \\ 0 & 0 & 4 \end{pmatrix} \begin{pmatrix} x_1 \\ x_2 \\ x_3 \end{pmatrix}= \begin{pmatrix} 3 \\ -1 \\ -4 \end{pmatrix}}
	\]
\psol{\itshape We can solve this using back-substitution:}
	\[
	\begin{aligned}
	\psol{0x_1 + 0x_2 + 4x_3= -4} &\psol{\quad\Longrightarrow\quad 4x_3= -4} &\psol{\quad\Longrightarrow\quad x_3= -1} \\
	\psol{0x_1 + 3x_2 + x_3= -1} &\psol{\quad\Longrightarrow\quad 3x_2 - 1= -1} &\psol{\quad\Longrightarrow\quad x_2= 0} \\
	\psol{1x_1 + 2x_2 - x_3= 3} &\psol{\quad\Longrightarrow\quad x_1 + 2(0) - (-1)= 3} &\psol{\quad\Longrightarrow\quad x_1= 2}
	\end{aligned}
	\]
\psol{\itshape Therefore, the solution is $\mathbf{x}= \begin{pmatrix} 2 \\ 0 \\ -1 \end{pmatrix}$.}
\end{enumerate}



% Question 2
\newpage
\question[10] Consider the invertible matrix $A= \begin{pmatrix} 1 & 3 & -2 \\ 2 & 6 & 5 \\ -3 & -1 & 1 \end{pmatrix}$. \par\vspace{0.3cm}

\begin{parts}
\part Explain why $A$ \textit{does not} have an LU decomposition without computing a PLU decomposition for $A$, i.e. do not use (b). \par\vspace{0.75cm}
\psol{\itshape Because $A$ is a square invertible matrix, $A$ has an LU if and only if all the leading principal minors are nonzero. The first leading principal minor is $M_1= \det(1)= 1 \neq 0$. However, $M_2= \det \begin{pmatrix} 1 & 3 \\ 2 & 6 \end{pmatrix}= 1(6) - 2(3)= 6 - 6= 0$. Therefore, $A$ does not have an LU decomposition. But $A$ must have a $PA= LU$ or $A= PLU$ decomposition.} \par\vspace{1.39cm}


\part Find a PLU decomposition for the matrix $A$. \pspace
	\begin{table}[!ht]
	\centering
	\begin{tabular}{C{3cm} C{3cm} C{3cm}} 
	\psol{Row Operations} & \psol{Current $L$} & \psol{Current $U$} \\ \clinevsol{\cline{1-3}}
	\\
	\begin{minipage}{3cm} \centering \psol{$-2R_1 + R_2 \to R_2$} \\ \psol{$3R_1 + R_3 \to R_3$} \end{minipage} & \psol{$\begin{pmatrix} 1 & 0 & 0 \\ ? & 1 & 0 \\ ? & ? & 1 \end{pmatrix}$} & \psol{$\begin{pmatrix} 1 & 3 & -2 \\ 2 & 6 & 5 \\ -3 & -1 & 1 \end{pmatrix}$} \\
	\\
	\psol{$R_2 \longleftrightarrow R_3$} & \psol{$\begin{pmatrix} 1 & 0 & 0 \\ 2 & 1 & 0 \\ -3 & ? & 1 \end{pmatrix}$} & \psol{$\begin{pmatrix} 1 & 3 & -2 \\ 0 & 0 & 9 \\ 0 & 8 & -5 \end{pmatrix}$} \\
	\\
	\psol{---} & \psol{$\underbrace{\begin{pmatrix} 1 & 0 & 0 \\ -3 & 1 & 0 \\ 2 & 0 & 1 \end{pmatrix}}_{L}$} & \psol{$\underbrace{\begin{pmatrix} 1 & 3 & -2 \\ 0 & 8 & -5 \\ 0 & 0 & 9 \end{pmatrix}}_{U}$}
	\end{tabular}
	\end{table} \par
\psol{\itshape We had to swap $R_2$ and $R_3$. The permutation that only swaps $R_2$ and $R_3$ for a $3 \times 3$ matrix is $P= \begin{pmatrix} 1 & 0 & 0 \\ 0 & 0 & 1 \\ 0 & 1 & 0 \end{pmatrix}$. We then know $PA= LU$. So, $A= P^{-1} LU$. But because $P$ is a permutation matrix, we know that $P^{-1}= P^T$. So, $A= P^T LU$. We have $P^T= \begin{pmatrix} 1 & 0 & 0 \\ 0 & 0 & 1 \\ 0 & 1 & 0 \end{pmatrix}$. Therefore, we have\dots}
	\[
	\psol{A= \underbrace{\begin{pmatrix} 1 & 0 & 0 \\ 0 & 0 & 1 \\ 0 & 1 & 0 \end{pmatrix}}_{P} \;\; \underbrace{\begin{pmatrix} 1 & 0 & 0 \\ -3 & 1 & 0 \\ 2 & 0 & 1 \end{pmatrix}}_{L} \;\; \underbrace{\begin{pmatrix} 1 & 3 & -2 \\ 0 & 8 & -5 \\ 0 & 0 & 9 \end{pmatrix}}_{U}}
	\]
\end{parts}



% Question 3
\newpage
\question[10] A matrix $A$ and its reduced-row echelon form, $R$, are given below.
	\[
	A= \begin{pmatrix} 1 & -2 & 0 & 0 & 3 \\ 3 & -6 & 1 & 2 & -1 \\ 2 & -4 & 1 & 0 & 6 \end{pmatrix}, \qquad  R= \begin{pmatrix} \vanbox{1} & -2 & 0 & 0 & 3 \\ 0 & 0 & \vanbox{1} & 0 & 0 \\ 0 & 0 & 0 & \vanbox{1} & -5 \end{pmatrix}
	\]

\begin{enumerate}[(a)]
\item Determine which $\mathbb{R}^n$ each of the spaces below are a subspace of and find their dimension. 
	\begin{table}[!ht]
	\centering
	\resizebox{0.7\columnwidth}{!}{%
	\begin{tabular}{|C{3cm}| C{3cm}| C{3cm}| C{3cm}|} \cline{1-4}
	& & & \\
	& $\operatorname{Col}(A)$ & $\operatorname{Row}(A)$ & $\operatorname{Null}(A)$ \\ 
	& & &\\ \cline{1-4}
	& & & \\ 
	Subspace of\dots & \psol{$\mathbb{R}^3$} & \psol{$\mathbb{R}^5$} & \psol{$\mathbb{R}^5$} \\ 
	& & & \\ \cline{1-4}
	& & & \\
	Dimension & \psol{$\rank A= 3$} & \psol{$\rank A= 3$} & \psol{$\nullity A= 2$} \\ 
	& & & \\ \cline{1-4}
	\end{tabular}%
	}
	\end{table}

\item Find a basis for $\operatorname{Col}(A)$. \pspace
\psol{\itshape We choose the pivot columns of $A$.}
	\[
	\resizebox{0.20\columnwidth}{!}{\psol{$\left\{ \begin{pmatrix} 1 \\ 3 \\ 2 \end{pmatrix}, \begin{pmatrix} 0 \\ 1 \\ 1 \end{pmatrix}, \begin{pmatrix} 0 \\ 2 \\ 0 \end{pmatrix} \right\}$}}
	\]

\item Find a basis for $\operatorname{Row}(A)$. \pspace
\psol{\itshape We choose the nonzero rows of $\text{rref}(A)$.}
	\[
	\resizebox{0.20\columnwidth}{!}{\psol{$\left\{ \begin{pmatrix} 1 \\ -2 \\ 0 \\ 0 \\ 3 \end{pmatrix}, \begin{pmatrix} 0 \\ 0 \\ 1 \\ 0 \\ 0 \end{pmatrix}, \begin{pmatrix} 0 \\ 0 \\ 0 \\ 1 \\ -5 \end{pmatrix} \right\}$}}
	\]

\item Find a basis for $\operatorname{Null}(A)$. \pspace
\psol{\itshape We know $\operatorname{Null}(A)$ is a subspace of $\mathbb{R}^5$, so we have five variables: $x_1, \ldots, x_5$. The basic variables are $x_1, x_3, x_4$ and the free variables are $x_2, x_5$. We know $x_1 - 2x_2 + 3x_5= 0$, $x_3= 0$, and $x_4 - 5x_5= 0$, i.e. $x_1= 2x_2 - 3x_5$, $x_3= 0$, and $x_4= 5x_5$. But then\dots}
	\[
	\resizebox{0.40\columnwidth}{!}{\psol{$\begin{pmatrix} x_1 \\ x_2 \\ x_3 \\ x_4 \\ x_5 \end{pmatrix}= \begin{pmatrix} 2x_2 - 3x_5 \\ x_2 \\ 0 \\ 5x_5 \\ x_5 \end{pmatrix}= x_2 \begin{pmatrix} 2 \\ 1 \\ 0 \\ 0 \\ 0 \end{pmatrix} + x_5 \begin{pmatrix} -3 \\ 0 \\ 0 \\ 5 \\ 1 \end{pmatrix}$}}
	\]
\psol{\itshape Therefore, a basis for the nullspace is\dots}
	\[
	\resizebox{0.15\columnwidth}{!}{\psol{$\left\{ \begin{pmatrix} 2 \\ 1 \\ 0 \\ 0 \\ 0 \end{pmatrix}, \begin{pmatrix} -3 \\ 0 \\ 0 \\ 5 \\ 1 \end{pmatrix} \right\}$}}
	\]
\end{enumerate}



% Question 4
\newpage
\question[10] Complete the parts given below.

\begin{enumerate}[(a)]
\item Compute the following \textit{without} explicitly computing the determinant: $\det \begin{pmatrix} 1 & 0 & 0 & 0 \\ 0 & 1 & 0 & 0 \\ -5 & 0 & 1 & 0 \\ 0 & 0 & 0 & 1 \end{pmatrix}$. \pspace
\psol{\itshape This is an elementary matrix that when multiplied on the left performs $-5R_1 + R_3 \to R_3$, which leaves the determinant unchanged. Therefore, we have\dots}
	\[
	\psol{\det \begin{pmatrix} 1 & 0 & 0 & 0 \\ 0 & 1 & 0 & 0 \\ -5 & 0 & 1 & 0 \\ 0 & 0 & 0 & 1 \end{pmatrix}= 1}
	\] \vfill

\item Find an elementary matrix that scales the second row of a $2 \times 2$ matrix by $\frac{1}{3}$. \pspace
	\[
	\psol{\begin{pmatrix} 1 & 0 \\ 0 & \frac{1}{3} \end{pmatrix}}
	\] \vfill

\item If $E= \begin{pmatrix} 1 & 0 & 0 \\ 0 & 3 & 0 \\ -2 & 0 & 1 \end{pmatrix}$ and $A$ is a $3 \times 3$ matrix, describe the matrix $EA$. \pspace
\psol{\itshape The product $EA$ has the same first row as $A$, the second row is $3$ times the second row of $A$, and the third row is the third row of $A$ replaced by twice the first row of $A$ subtracted from the third row of $A$. That is, $EA$ is $A$ after the elementary row operations $3R_2 \to R_2$ and $-2R_1 + R_3 \to R_3$.} \pspace \vfill


\item Compute the following \textit{without} explicitly multiplying the matrices:
	\[
	\begin{pmatrix} 0 & 0 & 1 & 0 \\ 0 & 0 & 0 & 1 \\ 1 & 0 & 0 & 0 \\ 0 & 1 & 0 & 0 \end{pmatrix} \begin{pmatrix} 11 & -15 & 18 & -20 \\ 8 & 5 & 16 & -9 \\ 0 & -3 & 10 & 1 \\ 3 & 17 & 0 & 4 \end{pmatrix}
	\] \pspace
\psol{\itshape The matrix on the left is a permutation matrix which swaps rows one and three, and swaps rows two and four, i.e. $R_1 \longleftrightarrow R_3$ and $R_2 \longleftrightarrow R_4$. Therefore, we have\dots}
	\[
	\psol{\begin{pmatrix} 0 & 0 & 1 & 0 \\ 0 & 0 & 0 & 1 \\ 1 & 0 & 0 & 0 \\ 0 & 1 & 0 & 0 \end{pmatrix} \begin{pmatrix} 11 & -15 & 18 & -20 \\ 8 & 5 & 16 & -9 \\ 0 & -3 & 10 & 1 \\ 3 & 17 & 0 & 4 \end{pmatrix}= \begin{pmatrix} 0 & -3 & 10 & 1 \\ 3 & 17 & 0 & 4 \\ 11 & -15 & 18 & -20 \\ 8 & 5 & 16 & -9 \end{pmatrix}}
	\] \vfill
\end{enumerate}



% Question 5
\newpage
\question[10] Complete the following parts:

\begin{enumerate}[(a)]
\item Is the set $\left\{ \begin{pmatrix} 1 \\ -2 \\ 0 \\ 5 \end{pmatrix}, \begin{pmatrix} 6 \\ -1 \\ 4 \\ 3 \end{pmatrix}, \begin{pmatrix} 0 \\ 0 \\ 0 \\ 0 \end{pmatrix}, \begin{pmatrix} 5 \\ -5 \\ 0 \\ 1 \end{pmatrix} \right\}$ linearly independent? Explain. \vfill
\psol{\itshape No, the set contains the zero vector, $\mathbf{0}$. No collection of vectors containing the zero vector can be linearly independent.} \vfill

\item Can a set of 13 vectors in $\mathbb{R}^5$ be linearly independent? Explain. \pspace
\psol{\itshape No, the vectors cannot be linearly independent. The span the set is clearly a subspace of $\mathbb{R}^5$. The only possible dimensions of a subspace of $\mathbb{R}^5$ are $0, 1, \ldots, 5$. If the given vectors were linearly independent, then they would form a basis for the subspace generated by their span. This subspace would then have dimension 13, which is not possible. Therefore, the vectors must be linearly independent. Alternatively, we know a subspace of dimension $D$ cannot have a set of $N > D$ linearly independent vectors.} \vfill

\item Find a basis for the subspace $W= \{ (x, y, z) \in \mathbb{R}^3 \colon x + 2y - z= 0 \}$ and find its dimension. \pspace
\psol{\itshape We know that $x + 2y - z= 0$ if and only if $z= x + 2y$. But then\dots}
	\[
	\psol{\begin{pmatrix} x \\ y \\ z \end{pmatrix}= \begin{pmatrix} x \\ y \\ x + 2y \end{pmatrix}= x \begin{pmatrix} 1 \\ 0 \\ 1 \end{pmatrix} + y \begin{pmatrix} 0 \\ 1 \\ 2 \end{pmatrix}}
	\]
\psol{\itshape Therefore, $W= \mspan \left\{ \begin{pmatrix} 1 \\ 0 \\ 1 \end{pmatrix}, \begin{pmatrix} 0 \\ 1 \\ 2 \end{pmatrix} \right\}$. But these two vectors are not parallel. So, this set is linearly independent. But then the set is a basis for $W$. This means that $\dim W= 2$. Indeed, $x + 2y - z= 0$ is the equation of a plane in $\mathbb{R}^3$ passing through the origin.} \vfill
\end{enumerate}



% Question 6
\newpage
\question[10] Mark the following statements True ($T$) or False ($F$). You \textit{do not} need to justify your answer. \par\vspace{1cm}

\begin{enumerate}[(a)]
\item \underline{\itshape\hspace{0.65cm}\psol{F}\hspace{0.65cm}}: If $\nullity(A)= 0$, i.e. $\dim \operatorname{null} A= \{ \mathbf{0} \}$, then $A\mathbf{x}= \mathbf{0}$ does not have a solution. \vfill % F

\item \underline{\itshape\hspace{0.65cm}\psol{F}\hspace{0.65cm}}: The number of independent rows and independent columns of a matrix can be different. \vfill % F

\item \underline{\itshape\hspace{0.6cm}\psol{T}\hspace{0.6cm}}: A $12 \times 9$ matrix $A$ can have rank at most $9$. \vfill % T

\item \underline{\itshape\hspace{0.6cm}\psol{T}\hspace{0.6cm}}: The columns of \textit{any} $6 \times 7$ matrix are linearly dependent. \vfill % T

\item \underline{\itshape\hspace{0.6cm}\psol{T}\hspace{0.6cm}}: If $A$ is a $3 \times 4$ matrix with rank 3, then $A\mathbf{x}= \mathbf{b}$ \textit{always} has a solution but the solution is \textit{never} unique. \par\vspace{1cm} % T
\end{enumerate}

\end{questions}
\end{document}